\SourceLeaf{026}
\Article{14}
\textsc{Lemma.} \textit{Si quantitas $r$ angulusque $\varphi$ ita sunt determinati, ut habeantur aequationes}
\[
\begin{aligned}
&r^m\cos m\varphi+Ar^{m-1}\cos(m-1)\varphi+Br^{m-2}\cos(m-2)\varphi+\text{etc.}\\
&\hspace{3em}+Krr\cos2\varphi+Lr\cos\varphi+M=0\quad.\quad.\quad[1]\\
&r^m\sin m\varphi+Ar^{m-1}\sin(m-1)\varphi+Br^{m-2}\sin(m-2)\varphi+\text{etc.}\\
&\hspace{3em}+Krr\sin2\varphi+Lr\sin\varphi=0\quad.\quad.\quad[2]
\end{aligned}
\]
\textit{functio $x^m+Ax^{m-1}+Bx^{m-2}+\text{etc.}+Kxx+Lx+M=X$ divisibilis erit per factorem duplicem $xx-2\cos\varphi\mathbin{.}rx+rr$, si modo $r\sin\varphi$ non $=0$; si vero $r\sin\varphi=0$, eadem functio divisibilis erit per factorem simplicem $x-r\cos\varphi$.}

\textit{Demonstr.} I. Ex art. praec. omnes sequentes quantitates divisibiles erunt per $xx-2\cos\varphi\mathbin{.}rx+rr$:
{\small
\[
\setlength{\arraycolsep}{2pt}
\begin{array}{rll}
\sin\varphi\mathbin{.}rx^m&-\sin m\varphi\mathbin{.}r^m x&+\sin(m-1)\varphi\mathbin{.}r^{m+1}\\
A\sin\varphi\mathbin{.}rx^{m-1}&-A\sin(m-1)\varphi\mathbin{.}r^{m-1}x&+A\sin(m-2)\varphi\mathbin{.}r^m\\
B\sin\varphi\mathbin{.}rx^{m-2}&-B\sin(m-2)\varphi\mathbin{.}r^{m-2}x&+B\sin(m-3)\varphi\mathbin{.}r^{m-1}\\
&\text{etc.}&\text{etc.}\\
K\sin\varphi\mathbin{.}rxx&-K\sin2\varphi\mathbin{.}rrx&+K\sin\varphi\mathbin{.}r^3\\
L\sin\varphi\mathbin{.}rx&-L\sin\varphi\mathbin{.}rx&\qquad *\\
M\sin\varphi\mathbin{.}r&\qquad *&+M\sin(-\varphi)\mathbin{.}r
\end{array}
\]
}
Quamobrem etiam summa harum quantitatum per $xx-2\cos\varphi\mathbin{.}rx+rr$ divisibilis erit. At singularum partes primae constituunt summam $\sin\varphi\mathbin{.}rX$; secundae additae dant $0$, propter $[2]$; tertiarum vero aggregatum quoque evanescere, facile perspicitur, si $[1]$ multiplicatur per $\sin\varphi$, $[2]$ per $\cos\varphi$, productumque illud ab hoc subducitur. Unde sequitur, functionem $\sin\varphi\mathbin{.}rX$ divisibilem esse per $xx-2\cos\varphi\mathbin{.}rx+rr$, adeoque, nisi fuerit $r\sin\varphi=0$, etiam functionem $X$. Q. E. P.

II. Si vero $r\sin\varphi=0$, erit aut $r=0$ aut $\sin\varphi=0$. In casu priori erit $M=0$, propter $[1]$, adeoque $X$ per $x$ sive per $x-r\cos\varphi$ divisibilis; in posteriori erit $\cos\varphi=\pm1$, $\cos2\varphi=+1$, $\cos3\varphi=\pm1$ et generaliter $\cos n\varphi=\cos\varphi^n$. Quare propter $[1]$ fiet $X=0$, statuendo $x=r\cos\varphi$, et proin functio $X$ per $x-r\cos\varphi$ erit divisibilis. Q. E. S.

\SourceLeaf{027}
\Article{15}
Theorema praecedens plerumque adiumento quantitatum imaginariarum demonstratur, vid. \textsc{Euler} \textit{Introd. in Anal. Inf.} T. I. p. 110; operae pretium esse duxi, ostendere, quomodo aeque facile absque illarum auxilio erui possit. Manifestum iam est, ad demonstrationem theorematis nostri nihil aliud requiri quam ut ostendatur: \textit{Proposita functione quacunque $X$ formae $x^m+Ax^{m-1}+Bx^{m-2}+\text{etc.}+Lx+M$, $r$ et $\varphi$ ita determinari posse, ut aequationes $[1]$ et $[2]$ locum habeant.} Hinc enim sequitur, $X$ habere factorem realem primi vel secundi gradus; divisio autem necessario producet quotientem realem inferioris gradus, qui ex eadem ratione quoque factorem primi vel secundi gradus habebit. Per continuationem huius operationis $X$ tandem in factores reales simplices vel duplices resolvetur. Illud itaque theorema demonstrare, propositum est sequentium disquisitionum.

\Article{16}
Concipiatur planum fixum infinitum (planum tabulae, fig. 1), et in hoc recta fixa infinita $GC$ per punctum fixum $C$ transiens. Assumta aliqua longitudine pro unitate ut omnes rectae per numeros exprimi possint, erigatur in quovis puncto plani $P$, cuius distantia a centro $C$ est $r$ angulusque $GCP=\varphi$, perpendiculum aequale valori expressionis
\[
r^m\sin m\varphi+Ar^{m-1}\sin(m-1)\varphi+\text{etc.}+Lr\sin\varphi
\]
quem brevitatis gratia in sequentibus semper per $T$ designabo. Distantiam $r$ semper tamquam positivam considero, et pro punctis, quae axi ab altera parte iacent, angulus $\varphi$ aut tamquam duobus rectis maior, aut tamquam negativus (quod hic eodem redit) spectari debet. Extremitates horum perpendiculorum (quae pro valore positivo ipsius $T$ supra planum accipiendae sunt, pro negativo infra, pro evanescente in plano ipso) erunt ad superficiem curvam continuam quaquaversum infinitam, quam brevitatis gratia in sequentibus \textit{superficiem primam} vocabo. Prorsus simili modo ad idem planum et centrum eundemque axem referatur alia superficies, cuius altitudo supra quodvis plani punctum sit
\[
r^m\cos m\varphi+Ar^{m-1}\cos(m-1)\varphi+\text{etc.}+Lr\cos\varphi+M
\]
quam expressionem brevitatis gratia semper per $U$ denotabo. Superficiem vero hanc, quae etiam continua et quaquaversum infinita erit, per denominationem
